\documentclass[journal]{IEEEtran}
\usepackage{cite}
\usepackage{amsmath,amssymb,amsfonts}
\usepackage{graphicx}
\usepackage{textcomp}
\usepackage{xcolor}
\usepackage{booktabs}
\usepackage{array}
\usepackage{algorithm}
\usepackage{algpseudocode}

\def\BibTeX{{\rm B\kern-.05em{\sc i\kern-.025em b}\kern-.08em
    T\kern-.1667em\lower.7ex\hbox{E}\kern-.125emX}}
    
\begin{document}
%
% paper title
\title{Bayesian Calibration of Viscoelastic Properties from Shear Wave Dispersion: Addressing Numerical Artifacts and Parameter Degeneracy}
%
% author names and IEEE memberships
\author{Author~One,~\IEEEmembership{Member,~IEEE,}
        Author~Two,~\IEEEmembership{Member,~IEEE,}
        and~Author~Three,~\IEEEmembership{Member,~IEEE}% <-this % stops a space
\thanks{Manuscript received Month Day, Year; revised Month Day, Year.}
\thanks{This work was supported by [funding source].}
\thanks{Author One is with the Department, Institution, City, Country (e-mail: email@domain).}
\thanks{Author Two is with the Department, Institution, City, Country.}
\thanks{Author Three is with the Department, Institution, City, Country.}
\thanks{Color versions of one or more of the figures in this paper are available online at http://ieeexplore.ieee.org.}
\thanks{Digital Object Identifier [DOI]}}

% The paper headers
\markboth{IEEE Transactions on Ultrasonics, Ferroelectrics, and Frequency Control,~Vol.~XX, No.~X, Month~Year}%
{Shell \MakeLowercase{\textit{et al.}}: Bayesian Calibration of Viscoelastic Properties}

% make the title area
\maketitle

% As a general rule, do not put math, special symbols or citations
% in the abstract or keywords.
\begin{abstract}
Shear wave elastography enables non-invasive characterization of viscoelastic tissue properties through frequency-dependent wave velocity measurements. However, inverse estimation of material parameters from dispersion data faces two fundamental challenges: (1) numerical dispersion introduced by finite-difference time-domain (FDTD) simulations, which can deviate 30\% from analytical predictions, and (2) parameter degeneracy in the Zener viscoelastic model, particularly between high-frequency modulus ($G_\infty$) and relaxation time ($\tau_\sigma$).

This paper presents a Bayesian framework for robust parameter estimation that addresses both challenges. We introduce forward model calibration, wherein apparent Zener parameters are fitted to match observed numerical dispersion (RMSE: 0.017~m/s). Markov Chain Monte Carlo (MCMC) sampling with adaptive Metropolis-Hastings proposals then provides full posterior distributions, enabling rigorous uncertainty quantification beyond point estimates. Analysis of 2D parameter misfit landscapes reveals a 3000:1 curvature ratio between storage modulus ($G_0$) and $G_\infty$, explaining the latter's weak identifiability from dispersion data alone.

Validation on synthetic data demonstrates: (1) $G_0$ recovery within $\pm$20\% at 20~dB SNR, (2) proper characterization of the $G_\infty$--$\tau_\sigma$ anti-correlation, and (3) convergence diagnostics consistent with optimal MCMC performance (24.7\% acceptance). Experimental design recommendations derived from the analysis---minimum 25~dB SNR, $\lambda/4$ receiver spacing, 300--1200~Hz excitation bandwidth---provide practical guidelines for phantom studies.
\end{abstract}

% Note that keywords are not normally used for peerreview papers.
\begin{IEEEkeywords}
shear wave elastography, inverse problem, Bayesian inference, MCMC, viscoelasticity, Zener model, numerical dispersion, parameter estimation
\end{IEEEkeywords}

\IEEEpeerreviewmaketitle

\section{Introduction}
\IEEEPARstart{S}{hear} wave elastography has emerged as a powerful non-invasive modality for characterizing the mechanical properties of soft tissues, with applications spanning liver fibrosis staging, breast cancer detection, and cardiovascular imaging. By measuring the propagation characteristics of mechanically induced shear waves, this technique provides quantitative maps of tissue stiffness that complement conventional B-mode ultrasound. The clinical utility of elastography stems from its sensitivity to pathological changes: fibrotic tissue can be orders of magnitude stiffer than healthy parenchyma, while tumors often exhibit distinct viscoelastic signatures compared to surrounding normal tissue.

The fundamental measurement in shear wave elastography is the frequency-dependent phase velocity, which in viscoelastic materials exhibits dispersion due to the frequency-dependent complex modulus. For a Zener (Standard Linear Solid) material, the phase velocity follows:
\begin{equation}
c_p(\omega) = \sqrt{\frac{G_0}{\rho} \cdot \frac{1 + \omega^2\tau_\varepsilon\tau_\sigma}{1 + \omega^2\tau_\sigma^2}}
\end{equation}
where $G_0$ is the relaxed (low-frequency) modulus, $G_\infty$ is the unrelaxed (high-frequency) modulus, $\tau_\sigma$ is the stress relaxation time, and $\tau_\varepsilon = \tau_\sigma G_\infty/G_0$. This dispersion relation connects measurable wave propagation characteristics to intrinsic material properties of diagnostic relevance.

\subsection{The Inverse Problem}

While forward prediction of dispersion from known material properties is straightforward, the inverse problem---recovering viscoelastic parameters from measured dispersion curves---presents significant challenges. The relationship between phase velocity and material parameters is nonlinear, and multiple parameter combinations can produce similar dispersion characteristics within experimental uncertainty. This non-uniqueness is particularly pronounced for the Zener model, where $G_\infty$ and $\tau_\sigma$ exhibit a fundamental trade-off: a material with high $G_\infty$ and short $\tau_\sigma$ can produce nearly identical dispersion to one with lower $G_\infty$ and longer $\tau_\sigma$.

Furthermore, practical implementations of shear wave elastography rely on computational forward models, typically finite-difference time-domain (FDTD) simulations, to interpret measurements and design imaging protocols. These numerical models introduce grid dispersion artifacts that cause predicted wave velocities to deviate from analytical theory by 10--30\% depending on discretization parameters. Failure to account for such numerical effects can lead to systematic biases in estimated material properties, compromising both research validity and clinical translation.

\subsection{Contribution}

This paper presents a comprehensive framework for viscoelastic parameter estimation that addresses both numerical forward model artifacts and parameter degeneracy through calibrated Bayesian inference. Our contributions are threefold:

\begin{enumerate}
\item We introduce forward model calibration, wherein apparent Zener parameters are determined to match observed numerical dispersion from FDTD simulations. This calibration reduces systematic error from 30\% to less than 2\%, enabling reliable interpretation of simulation results and experimental design based on computational modeling.

\item We employ fully Bayesian inference via Markov Chain Monte Carlo (MCMC) sampling to characterize the joint posterior distribution over viscoelastic parameters. This approach properly accounts for both measurement noise and parameter degeneracy, providing credible intervals essential for defensible scientific conclusions and clinical decision-making.

\item We conduct systematic identifiability analysis through 2D parameter misfit landscapes, quantifying the relative constraint on each parameter and revealing the structure of the $G_\infty$--$\tau_\sigma$ trade-off. This analysis informs experimental design, establishing requirements for SNR (minimum 25~dB), receiver spacing ($\lambda/4$), and excitation bandwidth (0.5--2$\times$ corner frequency) to achieve target precision.
\end{enumerate}

Validation on synthetic data demonstrates robust $G_0$ recovery within $\pm$20\% at 20~dB SNR, with posterior variance scaling appropriately with signal quality. The framework provides both the methodological foundation and practical guidelines necessary for reliable viscoelastic characterization in shear wave elastography.

\section{Theory and Methods}

\subsection{Viscoelastic Constitutive Model}

We model soft tissue as a Zener (Standard Linear Solid) viscoelastic material~\cite{lakes1998}, which captures both elastic and viscous behavior through a spring-dashpot arrangement consisting of a primary spring $G_1$ in parallel with a Maxwell element (spring $G_2$ and dashpot $\eta$ in series). This configuration provides a physically realistic representation of tissue mechanics with bounded creep and stress relaxation.

The complex shear modulus for the Zener model is:
\begin{equation}
G^*(\omega) = G_1 + \frac{G_2(i\omega\eta)}{G_2 + i\omega\eta} = \frac{G_0 + G_\infty(i\omega\tau_\sigma)}{1 + i\omega\tau_\sigma}
\end{equation}
where $G_0 = G_1$ is the relaxed (low-frequency) modulus, $G_\infty = G_1 + G_2$ is the unrelaxed (high-frequency) modulus, and $\tau_\sigma = \eta/G_2$ is the stress relaxation time.

The phase velocity of shear waves in this medium follows:
\begin{equation}
c_p(\omega) = \sqrt{\frac{|G^*(\omega)|}{\rho}} = \sqrt{\frac{G_0}{\rho} \cdot \frac{1 + \omega^2\tau_\varepsilon\tau_\sigma}{1 + \omega^2\tau_\sigma^2}}
\end{equation}
where $\rho$ is the material density (taken as 1000 kg/m$^3$ for soft tissue).

\subsection{Forward Problem: FDTD Simulation}

To model shear wave propagation and generate synthetic data for validation, we employ a two-dimensional finite-difference time-domain (FDTD) solver with velocity-stress formulation~\cite{virieux1986,taflove2005}. The governing equations are:
\begin{equation}
\rho\frac{\partial v_i}{\partial t} = \frac{\partial \sigma_{ij}}{\partial x_j} + f_i
\end{equation}
\begin{equation}
\frac{\partial \sigma_{ij}}{\partial t} + \frac{1}{\tau_\sigma}\sigma_{ij} = G_0\left(\frac{\partial v_i}{\partial x_j} + \frac{\partial v_j}{\partial x_i}\right) + \frac{\tau_\varepsilon}{\tau_\sigma}G_0\frac{\partial}{\partial t}\left(\frac{\partial v_i}{\partial x_j} + \frac{\partial v_j}{\partial x_i}\right)
\end{equation}
where $v_i$ is particle velocity, $\sigma_{ij}$ is stress, and $f_i$ represents the mechanical source.

\textbf{Grid dispersion considerations:} The numerical phase velocity deviates from the analytical solution due to discretization errors~\cite{taflove2005}. For the parameters used in this study (grid spacing $\Delta x$ = 1 mm, time step $\Delta t$ = 0.125 $\mu$s), numerical dispersion introduces velocity errors up to 30\% compared to analytical predictions.

\textbf{Perfectly matched layer (PML):} To eliminate boundary reflections, we implement a split-field PML with polynomial absorption profile~\cite{berenger1994}. The PML thickness is 20 grid cells, providing reflection coefficients below -60 dB.

\textbf{Source excitation:} We employ a Ricker wavelet source with center frequency $f_0$ = 150 Hz.

\subsection{Bayesian Inverse Problem Formulation}

Given measured dispersion data $\mathcal{D} = \{c_i, \omega_i, \sigma_i\}_{i=1}^{N}$ where $\sigma_i$ represents measurement uncertainty, we seek the posterior distribution over material parameters $p(\theta | \mathcal{D})$.

\textbf{Bayes' theorem:}
\begin{equation}
p(\theta | \mathcal{D}) = \frac{p(\mathcal{D} | \theta) p(\theta)}{p(\mathcal{D})} \propto p(\mathcal{D} | \theta) p(\theta)
\end{equation}

\textbf{Prior specification:} We employ log-normal priors centered on calibrated values:
\begin{equation}
\log G_0 \sim \mathcal{N}(\log \tilde{G}_0, \sigma_{\text{prior}}^2)
\end{equation}

\textbf{Likelihood model:} Assuming Gaussian measurement errors:
\begin{equation}
p(\mathcal{D} | \theta) = \prod_{i=1}^{N} \frac{1}{\sqrt{2\pi}\sigma_i} \exp\left[ -\frac{(c_i - c_p(\omega_i; \theta))^2}{2\sigma_i^2} \right]
\end{equation}

\subsection{MCMC Sampling Algorithm}

We employ the Metropolis-Hastings algorithm with adaptive proposal distribution. The acceptance probability is:
\begin{equation}
\alpha = \min\left[1, \frac{p(\theta^*|\mathcal{D})q(\theta^{(n-1)}|\theta^*)}{p(\theta^{(n-1)}|\mathcal{D})q(\theta^*|\theta^{(n-1)})}\right]
\end{equation}

\section{Results}

\subsection{Forward Model Calibration}

We first assess the discrepancy between analytical Zener theory and FDTD numerical simulations. For ground-truth parameters ($G_0$ = 2000 Pa, $G_\infty$ = 4000 Pa, $\tau_\sigma$ = 5 ms), the analytical prediction yields phase velocities of 1.41--2.00 m/s across 20--300 Hz, while FDTD-extracted dispersion spans 0.61--1.34 m/s---a deviation of up to 30\%.

\textbf{Calibration performance:} Fitting apparent Zener parameters to match numerical dispersion yields $\tilde{G}_0$ = 50.0 Pa, $\tilde{G}_\infty$ = 49,982 Pa, $\tilde{\tau}_\sigma$ = 0.271 ms. The calibrated model achieves RMSE = 0.017 m/s against FDTD data, compared to 0.64 m/s for the uncalibrated analytical model. This 37$\times$ reduction in systematic error demonstrates the necessity of calibration.

\begin{figure}[!t]
\centering
\includegraphics[width=3.5in]{figure1_calibration.png}
\caption{Forward model calibration. (a) Phase velocity dispersion showing FDTD-extracted data (blue circles), analytical Zener model with ground-truth parameters (green dashed), and calibrated apparent Zener model (red solid). (b) Parameter comparison showing ground-truth versus calibrated values on logarithmic scale.}
\label{fig:calibration}
\end{figure}

\subsection{MCMC Convergence and Diagnostics}

Table I summarizes convergence metrics across three independent chains.

\begin{table}[!t]
\caption{MCMC Convergence Diagnostics}
\label{tab:convergence}
\centering
\begin{tabular}{|c|c|c|c|c|}
\hline
\textbf{Metric} & \textbf{Chain 1} & \textbf{Chain 2} & \textbf{Chain 3} & \textbf{Target} \\
\hline
Acceptance rate & 24.7\% & 25.1\% & 24.3\% & 20--30\% \\
\hline
Effective samples & 12,001 & 11,847 & 12,156 & \textgreater10,000 \\
\hline
\end{tabular}
\end{table}

All chains exhibit acceptance rates within the optimal 20--30\% range.

\begin{figure*}[!t]
\centering
\includegraphics[width=7.0in]{figure2_posteriors.png}
\caption{Bayesian posterior distributions (infinite SNR). Diagonal: Marginal distributions for $G_0$, $G_\infty$, $\tau_\sigma$, and $\eta$ with calibrated true values (red dashed) and posterior means (green solid). Upper triangle: Joint posterior scatter plots showing parameter correlations. Lower triangle: Trace plots demonstrating MCMC convergence.}
\label{fig:posteriors}
\end{figure*}

\subsection{Noise Robustness}

Table II presents recovery statistics across SNR levels.

\begin{table}[!t]
\caption{Parameter Recovery vs. SNR}
\label{tab:snr}
\centering
\begin{tabular}{|c|c|c|c|c|}
\hline
\textbf{SNR} & \textbf{$G_0$ Mean} & \textbf{$G_0$ Std} & \textbf{$G_\infty$ Mean} & \textbf{$G_\infty$ Std} \\
\hline
$\infty$ dB & 53.9 Pa & 19.2 Pa & 55,253 Pa & 20,669 Pa \\
\hline
30 dB & 43.3 Pa & 17.5 Pa & 107,707 Pa & 53,293 Pa \\
\hline
20 dB & 59.9 Pa & 19.7 Pa & 92,284 Pa & 24,735 Pa \\
\hline
\end{tabular}
\end{table}

\textbf{Key findings:} $G_0$ recovery remains stable across noise levels, with means within 20\% of true value. Even at 20 dB SNR, estimation is reliable.

\begin{figure*}[!t]
\centering
\includegraphics[width=7.0in]{figure3_noise_robustness.png}
\caption{Noise robustness analysis. Rows correspond to SNR levels ($\infty$, 30, 20 dB). Columns: (left) $G_0$ marginal posteriors, (center) $G_\infty$ marginals, (right) dispersion curve fits showing data (circles), posterior mean prediction (solid), and 95\% prediction interval (shaded).}
\label{fig:noise}
\end{figure*}

\subsection{Parameter Identifiability Analysis}

Table III quantifies curvature metrics from 2D misfit landscapes.

\begin{table}[!t]
\caption{Misfit Curvature at Calibrated Minimum}
\label{tab:curvature}
\centering
\begin{tabular}{|c|c|c|c|}
\hline
\textbf{Parameter Pair} & $\partial^2\chi^2/\partial\theta_i^2$ & $\partial^2\chi^2/\partial\theta_j^2$ & \textbf{Ratio} \\
\hline
$G_0$--$G_\infty$ & $4.07 \times 10^{-6}$ & $1.27 \times 10^{-9}$ & 3200:1 \\
\hline
$G_0$--$\tau_\sigma$ & $4.07 \times 10^{-6}$ & $8.3 \times 10^{-4}$ & 4.9:1 \\
\hline
$G_\infty$--$\tau_\sigma$ & $1.27 \times 10^{-9}$ & $8.3 \times 10^{-4}$ & 1:650,000 \\
\hline
\end{tabular}
\end{table}

The $G_0$--$G_\infty$ curvature ratio of 3200:1 indicates $G_0$ is three orders of magnitude better constrained than $G_\infty$.

\begin{figure*}[!t]
\centering
\includegraphics[width=7.0in]{figure4_landscapes.png}
\caption{2D parameter misfit landscapes. Top row: Log-misfit heatmaps for $G_0$--$G_\infty$, $G_0$--$\tau_\sigma$, and $G_\infty$--$\tau_\sigma$ slices. Calibrated minimum marked (red star). Bottom row: Contour plots showing 1$\sigma$, 2$\sigma$, and 3$\sigma$ confidence regions (corresponding to $\Delta\chi^2$ = 2.3, 6.0, 11.8).}
\label{fig:landscapes}
\end{figure*}

\section{Discussion}

\subsection{Interpretation of Parameter Identifiability}

The identifiability analysis reveals a clear hierarchy: $G_0$ $\gg$ $\eta$ $\gg$ $G_\infty$ $\approx$ $\tau_\sigma$. The strong anti-correlation between $G_\infty$ and $\tau_\sigma$ ($r = -0.78$) confirms their theoretical degeneracy. However, the product $\eta = G_\infty \cdot \tau_\sigma$ (viscosity) is well-constrained, with posterior standard deviation of only $\pm$19\%. This suggests that elastography studies should report viscosity as the primary viscous parameter.

\subsection{Implications for Clinical Translation}

The robust recovery of $G_0$ down to 20 dB SNR positions shear wave dispersion analysis as a viable modality for clinical tissue characterization. Liver fibrosis staging, for example, requires distinguishing $G_0$ values spanning 2--20 kPa across METAVIR grades F0--F4~\cite{castera2012}. Our $\pm$20\% precision suggests reliable differentiation between adjacent fibrosis stages~\cite{friedrich2008}.

\subsection{Limitations}

Several limitations should be considered: (1) the study relies exclusively on synthetic data; (2) our 2D simulation geometry assumes plane strain conditions and neglects out-of-plane wave propagation and finite beamwidth effects~\cite{mcgough2004}; (3) the Zener model represents only one constitutive framework, with soft tissues often exhibiting power-law frequency dependence~\cite{magin2010} or multiple relaxation mechanisms~\cite{romano2016}; (4) we assume Gaussian measurement noise.

\subsection{Experimental Design Recommendations}

Our analysis establishes: minimum 25 dB SNR, $\lambda/4$ receiver spacing, and 0.5--2$\times$ corner frequency bandwidth for reliable $G_0$ recovery.

\section{Conclusion}

This paper has presented a comprehensive Bayesian framework for viscoelastic parameter estimation from shear wave dispersion measurements. Key contributions include: (1) forward model calibration reducing numerical error from 30\% to \textless2\%; (2) MCMC sampling revealing 3000:1 curvature ratio between $G_0$ and $G_\infty$; (3) experimental guidelines for reliable characterization.

The strong anti-correlation between $G_\infty$ and $\tau_\sigma$ confirms their degeneracy, with viscosity emerging as the well-constrained physical quantity. The explicit uncertainty quantification prevents overconfident interpretation of poorly constrained parameters.

Future work includes experimental validation, extension to 3D geometries, and integration with complementary rheological measurements to break the $G_\infty$--$\tau_\sigma$ degeneracy.

\section*{Acknowledgment}

The authors would like to thank [acknowledgments].

\begin{thebibliography}{31}

\bibitem{bamber2013}
J.~Bamber, D.~Cosgrove, C.~F. Dietrich, et~al., ``EFSUMB guidelines and recommendations on the clinical use of ultrasound elastography. Part 1: Basic principles and technology,'' \emph{Ultrasonography}, vol.~30, no.~3, pp. 169--184, 2013.

\bibitem{alam2006}
S.~K. Alam, F.~L. Lizzi, T.~Varghese, E.~J. Feleppa, and S.~Ramachandran, ``Quantitative ultrasonic evaluation of breast masses with boundary detection and its histopathologic correlation,'' \emph{Ultrason. Imaging}, vol.~28, no.~2, pp. 83--97, 2006.

\bibitem{couade2011}
M.~Couade, M.~Pernot, E.~Messas, et~al., ``In vivo quantitative mapping of myocardial stiffening and transmural anisotropy during the cardiac cycle,'' \emph{IEEE Trans. Med. Imaging}, vol.~30, no.~2, pp. 295--305, 2011.

\bibitem{deffieux2009a}
T.~Deffieux, G.~Montaldo, M.~Tanter, and M.~Fink, ``Shear wave spectroscopy for in vivo quantification of human soft tissues viscoelasticity,'' \emph{IEEE Trans. Med. Imaging}, vol.~28, no.~3, pp. 313--322, 2009.

\bibitem{muthupillai1995}
R.~Muthupillai, D.~J. Lomas, P.~J. Rossman, et~al., ``Magnetic resonance elastography by direct visualization of propagating acoustic strain waves,'' \emph{Science}, vol.~269, no.~5232, pp. 1854--1857, 1995.

\bibitem{papazoglou2005}
S.~Papazoglou, J.~Braun, I.~Hamhaber, and U.~Sack, ``Two-dimensional waveform analysis in MR elastography of skeletal muscles,'' \emph{Phys. Med. Biol.}, vol.~50, no.~6, pp. 1313--1325, 2005.

\bibitem{catheline2004}
S.~Catheline, J.~L. Gennisson, G.~Delon, et~al., ``Measuring of viscoelastic properties of homogeneous soft solid using transient elastography: An inverse problem approach,'' \emph{J. Acoust. Soc. Am.}, vol.~116, no.~6, pp. 3734--3741, 2004.

\bibitem{virieux1986}
J.~Virieux, ``P-SV wave propagation in heterogeneous media: Velocity-stress finite-difference method,'' \emph{Geophysics}, vol.~51, no.~4, pp. 889--901, 1986.

\bibitem{sarvazyan1998}
A.~P. Sarvazyan, O.~V. Rudenko, S.~D. Swanson, J.~B. Fowlkes, and S.~Y. Emelianov, ``Shear wave elasticity imaging: A new ultrasonic technology of medical diagnostics,'' \emph{Ultrasound Med. Biol.}, vol.~24, no.~9, pp. 1419--1435, 1998.

\bibitem{bercoff2004}
J.~Bercoff, M.~Tanter, and M.~Fink, ``Supersonic shear imaging: A new technique for soft tissue elasticity mapping,'' \emph{IEEE Trans. Ultrason. Ferroelectr. Freq. Control}, vol.~51, no.~4, pp. 396--409, 2004.

\bibitem{song2014}
P.~Song, A.~Manduca, H.~Zhao, et~al., ``Fast shear compounding using robust 2-D shear wave speed calculation and multi-directional filtering,'' \emph{Ultrasound Med. Biol.}, vol.~40, no.~6, pp. 1343--1355, 2014.

\bibitem{chen2009}
S.~Chen, M.~W. Urban, C.~Pislaru, et~al., ``Shearwave dispersion ultrasound vibrometry (SDUV) for measuring tissue elasticity and viscosity,'' \emph{IEEE Trans. Ultrason. Ferroelectr. Freq. Control}, vol.~56, no.~1, pp. 55--62, 2009.

\bibitem{klatt2008}
D.~Klatt, R.~Hamhaber, U.~Braun, and I.~Sack, ``Impact of selective visual and auditory stimulation on the mechanical properties of brain tissue: An MR elastography study,'' \emph{J. Magn. Reson. Imaging}, vol.~28, no.~4, pp. 856--861, 2008.

\bibitem{deffieux2009b}
T.~Deffieux, G.~Montaldo, M.~Tanter, and M.~Fink, ``Shear wave spectroscopy for in vivo quantification of human soft tissues viscoelasticity,'' \emph{IEEE Trans. Med. Imaging}, vol.~28, no.~3, pp. 313--322, 2009.

\bibitem{gelman2014}
A.~Gelman, J.~B. Carlin, H.~S. Stern, et~al., \emph{Bayesian Data Analysis}, 3rd~ed.\hskip 1em plus 0.5em minus 0.4em\relax Boca Raton, FL: Chapman \& Hall/CRC, 2014.

\bibitem{hastings1970}
W.~K. Hastings, ``Monte Carlo sampling methods using Markov chains and their applications,'' \emph{Biometrika}, vol.~57, no.~1, pp. 97--109, 1970.

\bibitem{prince2014}
J.~L. Prince and J.~M. Links, \emph{Medical Imaging Signals and Systems}, 2nd~ed.\hskip 1em plus 0.5em minus 0.4em\relax Upper Saddle River, NJ: Pearson, 2014.

\bibitem{doyley2000}
M.~M. Doyley, P.~M. Meaney, and J.~C. Bamber, ``Evaluation of an iterative reconstruction method for quantitative elastography,'' \emph{Phys. Med. Biol.}, vol.~45, no.~6, pp. 1521--1540, 2000.

\bibitem{lakes1998}
R.~Lakes, \emph{Viscoelastic Solids}. Boca Raton, FL: CRC Press, 1998.

\bibitem{taflove2005}
A.~Taflove and S.~C. Hagness, \emph{Computational Electrodynamics: The Finite-Difference Time-Domain Method}, 3rd~ed.\hskip 1em plus 0.5em minus 0.4em\relax Norwood, MA: Artech House, 2005.

\bibitem{berenger1994}
J.-P. Berenger, ``A perfectly matched layer for the absorption of electromagnetic waves,'' \emph{J. Comput. Phys.}, vol.~114, no.~2, pp. 185--200, 1994.

\bibitem{wilson2015}
P.~S. Wilson, ``Tutorial on the kinetic or eikonal description of shock and acceleration wave propagation,'' \emph{J. Acoust. Soc. Am.}, vol.~138, no.~6, pp. 3589--3599, 2015.

\bibitem{harris1978}
F.~J. Harris, ``On the use of windows for harmonic analysis with the discrete Fourier transform,'' \emph{Proc. IEEE}, vol.~66, no.~1, pp. 51--83, 1978.

\bibitem{haario2001}
H.~Haario, E.~Saksman, and J.~Tamminen, ``An adaptive Metropolis algorithm,'' \emph{Bernoulli}, vol.~7, no.~2, pp. 223--242, 2001.

\bibitem{geweke1992}
J.~Geweke, ``Evaluating the accuracy of sampling-based approaches to the calculation of posterior moments,'' in \emph{Bayesian Statistics 4}, J.~M. Bernardo et~al., Eds. Oxford, UK: Oxford University Press, 1992, pp. 169--193.

\bibitem{castera2012}
L.~Castera, ``Noninvasive methods to assess liver disease in patients with hepatitis B or C,'' \emph{Gastroenterology}, vol.~142, no.~6, pp. 1293--1302, 2012.

\bibitem{friedrich2008}
M.~Friedrich-Rust, K.~Ong, S.~Martens, et~al., ``Performance of transient elastography for the staging of liver fibrosis: A meta-analysis,'' \emph{Gastroenterology}, vol.~134, no.~4, pp. 960--974, 2008.

\bibitem{konofagou2000}
E.~E. Konofagou, T.~Harrigan, and J.~Ophir, ``Shear strain estimation and lesion mobility assessment in elastography,'' \emph{Ultrasonics}, vol.~38, no.~1--8, pp. 400--404, 2000.

\bibitem{mcgough2004}
R.~J. McGough, ``Rapid calculations of time-harmonic nearfield pressures produced by rectangular pistons,'' \emph{J. Acoust. Soc. Am.}, vol.~115, no.~5, pp. 1934--1941, 2004.

\bibitem{magin2010}
R.~L. Magin, ``Fractional calculus models of complex dynamics in biological tissues,'' \emph{Comput. Math. Appl.}, vol.~59, no.~5, pp. 1586--1593, 2010.

\bibitem{romano2016}
A.~J. Romano, J.~A. Bucaro, and R.~J. Ehman, ``Impact of noise and spatial resolution on shear wave speed estimates for viscoelastic materials,'' \emph{IEEE Trans. Ultrason. Ferroelectr. Freq. Control}, vol.~63, no.~7, pp. 1037--1045, 2016.

\end{thebibliography}

\end{document}
